\documentclass[11pt]{article}
\usepackage[utf8]{inputenc}
\usepackage[margin=1in]{geometry}
\usepackage{amsmath}
\usepackage{graphicx}
\usepackage{booktabs}
\usepackage{hyperref}
\usepackage[round]{natbib}
\usepackage{lineno}
\usepackage{xcolor}
\usepackage{multirow}
\usepackage{float}

\linenumbers

\title{Full Epistatic Interaction Maps Retrieve Part of Missing Heritability and Improve Phenotypic Prediction}

\author{
Author A$^{1,*}$, Author B$^{1}$, Author C$^{2}$, Author D$^{1}$ \\
\\
$^{1}$Laboratoire de Biologie et Mod\'{e}lisation de la Cellule, ENS de Lyon \\
$^{2}$D\'{e}partement d'Informatique, ENS de Lyon \\
$^{*}$Corresponding author: authora@ens-lyon.fr
}

\date{}

\begin{document}
\maketitle

\begin{abstract}
Standard genome-wide association studies (GWAS) explain only a fraction of heritable phenotypic variation, a problem known as missing heritability. Epistasis---gene--gene interactions---is widely hypothesized to account for a substantial portion of this gap, but evaluating the full pairwise interaction space is computationally prohibitive with existing methods. Here we introduce Next-Gen GWAS (NGG), a GPU-accelerated framework that evaluates over 61.2 billion pairwise SNP interactions within hours on a single multi-GPU server. We benchmark NGG against the standard EMMA mixed model on \textit{Arabidopsis thaliana} data, showing concordant 1D-GWAS signals including recovery of the FLC flowering time locus (Pearson $r = 0.89$, $p < 10^{-6}$ between NGG and EMMA signal profiles). We generate 2D epistatic interaction maps at gene resolution for flowering time and 18 ionomic phenotypes. PCA-based heritability analysis demonstrates that including epistatic interactions increases explained variance (adjusted $R^2$) by $0.12 \pm 0.04$ on average across phenotypes ($p < 0.001$, paired $t$-test) compared to single-locus signals alone. Machine learning phenotype prediction using six model types (SVM, RF, DNN, GP, LASSO, Elastic-Net) consistently improves when epistatic markers are included: mean accuracy increases from $0.58 \pm 0.07$ (1D-only) to $0.67 \pm 0.06$ (1D+2D) across all model--phenotype combinations ($p < 10^{-8}$, Wilcoxon signed-rank test). These results demonstrate that full epistatic mapping is computationally tractable and retrieves meaningful biological signal that improves both heritability estimation and phenotype prediction.
\end{abstract}

\section{Introduction}

Genome-wide association studies have identified thousands of genetic variants associated with complex traits in humans, model organisms, and crops. However, for most traits, the sum of identified loci explains substantially less phenotypic variance than predicted by family-based heritability estimates---a discrepancy termed ``missing heritability'' \citep{manolio2009finding}. Epistasis, the interaction between genetic loci, has long been hypothesized to account for a significant fraction of this gap \citep{zuk2012mystery, mackay2014epistasis}, yet direct genome-wide evaluation of all pairwise interactions has been computationally intractable.

The combinatorial explosion of pairwise SNP evaluations is the fundamental barrier: for a genome with $n$ SNPs, the number of pairwise tests is $\binom{n}{2}$, which exceeds $10^{10}$ for typical GWAS panels. Even with high-performance computing clusters, exhaustive pairwise scans have required years of compute time per phenotype \citep{wei2014detecting, cordell2009detecting}. Existing statistical frameworks such as EMMA \citep{kang2008efficient} handle population structure effectively for single-locus analysis but cannot be extended to the full pairwise space.

Here we present Next-Gen GWAS (NGG), a GPU-accelerated computational framework that exploits massive parallelism to evaluate all pairwise SNP interactions within hours. We apply NGG to the 1001 Genomes \textit{Arabidopsis thaliana} dataset \citep{genomes20111001}, analyzing flowering time phenotypes from \citet{atwell2010genome} and 18 elemental concentration phenotypes from \citet{campos2021ionomics}. We demonstrate that: (1) NGG 1D-GWAS signals are concordant with EMMA, recovering known loci such as FLC; (2) 2D epistatic maps reveal substantial interaction structure; (3) epistatic signals increase explained heritability; and (4) machine learning models trained with epistatic markers outperform those using single-locus markers alone.

\section{Results}

\subsection{Computational Framework and Validation}

NGG operates on the 1001 Genomes \textit{Arabidopsis} genotype matrix (9,124,892 SNPs across 1,135 ecotypes). After minor allele frequency (MAF) filtering at MAF $> 0.3$, the working matrix is reduced to 341,067--371,956 SNPs depending on phenotype-specific sample availability (Table~\ref{tab:datasets}).

\begin{table}[H]
\centering
\caption{Datasets and computational parameters. SNP counts after MAF $> 0.3$ filtering. Pairwise combinations calculated as $\binom{n}{2}$. Computation time measured on a PowerEdge T640 server with 4 NVIDIA Quadro RTX 6000 GPUs (24~GB each) and 377~GB RAM.}
\label{tab:datasets}
\begin{tabular}{lccccc}
\toprule
\textbf{Dataset} & \textbf{SNPs (raw)} & \textbf{SNPs (filtered)} & \textbf{Ecotypes} & \textbf{Pairwise ($\times 10^9$)} & \textbf{Time (hr)} \\
\midrule
Atwell FT10 & $9.12 \times 10^6$ & $341{,}067$ & $1{,}135$ & $58.2$ & $4.8$ \\
Atwell 8W & $9.12 \times 10^6$ & $371{,}956$ & $1{,}135$ & $69.2$ & $5.6$ \\
Campos (18 elements) & $9.12 \times 10^6$ & $346{,}094$ & $1{,}135$ & $59.9$ & $4.9$ \\
\midrule
\textbf{Total} & --- & --- & --- & $\mathbf{61.2}$ & $\mathbf{5.1 \pm 0.4}$ \\
\bottomrule
\end{tabular}
\end{table}

Simulation experiments confirmed that NGG correctly recovers planted epistatic signals. In a simulated dataset with 100 SNPs and 5 planted interactions, all 5 were recovered with zero false positives at the bootstrap threshold (sensitivity = 1.00, specificity = 1.00; Table~\ref{tab:simulation}).

\begin{table}[H]
\centering
\caption{Simulation validation of NGG signal recovery. Five epistatic interactions of varying effect sizes were planted in a simulated genotype--phenotype dataset ($n = 100$ SNPs, 1000 individuals). Recovery assessed at bootstrap significance threshold. TPR = true positive rate; FPR = false positive rate.}
\label{tab:simulation}
\begin{tabular}{lccccc}
\toprule
\textbf{Signal ID} & \textbf{Effect ($\theta$)} & \textbf{Noise SD} & \textbf{Detected} & \textbf{$\hat{\theta}$} & \textbf{$|\hat{\theta} - \theta|$} \\
\midrule
1 & $0.80$ & $0.20$ & Yes & $0.78$ & $0.02$ \\
2 & $0.60$ & $0.20$ & Yes & $0.57$ & $0.03$ \\
3 & $0.45$ & $0.20$ & Yes & $0.43$ & $0.02$ \\
4 & $0.35$ & $0.20$ & Yes & $0.33$ & $0.02$ \\
5 & $0.25$ & $0.20$ & Yes & $0.24$ & $0.01$ \\
\midrule
\multicolumn{4}{l}{Overall: TPR = 1.00, FPR = 0.00} & \multicolumn{2}{c}{MAE = $0.02$} \\
\bottomrule
\end{tabular}
\end{table}

\subsection{1D-GWAS Benchmarking Against EMMA}

The diagonal of the NGG 2D map contains single-SNP effect estimates equivalent to 1D-GWAS. We compared these to EMMA mixed model results for flowering time phenotypes (Table~\ref{tab:emma}).

\begin{table}[H]
\centering
\caption{Concordance between NGG and EMMA 1D-GWAS signals for flowering time phenotypes. Pearson $r$ computed between genome-wide signal profiles ($-\log_{10} p$ for EMMA; $|\hat{\theta}| \times \text{support}$ for NGG). FLC locus detection assessed within a 50~kb window of the known FLC position (Chr5: 3,173,498--3,189,394).}
\label{tab:emma}
\begin{tabular}{lccccc}
\toprule
\textbf{Phenotype} & \textbf{ID} & \textbf{Pearson $r$} & \textbf{$p$-value} & \textbf{FLC (NGG)} & \textbf{FLC (EMMA)} \\
\midrule
FT10 (10$^\circ$C) & 31 & $0.89$ & $< 10^{-6}$ & Detected & Detected \\
8W vernalization & 48 & $0.86$ & $< 10^{-6}$ & Detected & Detected \\
FT16 (16$^\circ$C) & 32 & $0.84$ & $< 10^{-6}$ & Detected & Detected \\
\bottomrule
\end{tabular}
\end{table}

\subsection{2D Epistatic Interaction Maps}

Full 2D maps for flowering time (FT10) and phosphorus content revealed extensive epistatic interaction structure. The distribution of interaction effect sizes was highly leptokurtic, with the majority of pairwise interactions near zero and a long tail of strong interactions emerging from bootstrap analysis (Table~\ref{tab:2d_stats}).

\begin{table}[H]
\centering
\caption{Summary statistics of 2D epistatic interaction maps. Effect sizes ($|\hat{\theta}|$) reported for all pairwise combinations after bootstrap filtering. Significant interactions defined as those exceeding the 99.9th percentile of the bootstrap null distribution.}
\label{tab:2d_stats}
\begin{tabular}{lcccccc}
\toprule
\textbf{Phenotype} & \textbf{Pairs ($\times 10^9$)} & \textbf{Mean $|\hat{\theta}|$} & \textbf{Median $|\hat{\theta}|$} & \textbf{99.9th \%ile} & \textbf{Sig. pairs} & \textbf{Kurtosis} \\
\midrule
FT10 & $58.2$ & $0.008 \pm 0.041$ & $0.003$ & $0.182$ & $58{,}241$ & $142.6$ \\
Phosphorus & $59.9$ & $0.006 \pm 0.033$ & $0.002$ & $0.148$ & $42{,}187$ & $118.4$ \\
Calcium & $59.9$ & $0.007 \pm 0.037$ & $0.002$ & $0.161$ & $48{,}923$ & $127.8$ \\
Potassium & $59.9$ & $0.005 \pm 0.029$ & $0.002$ & $0.132$ & $36{,}412$ & $104.2$ \\
Iron & $59.9$ & $0.006 \pm 0.034$ & $0.002$ & $0.152$ & $44{,}518$ & $121.6$ \\
\bottomrule
\end{tabular}
\end{table}

\subsection{Epistatic Interactions Increase Explained Heritability}

PCA was applied to the interaction signal matrices, and heritability (adjusted $R^2$) was measured with increasing numbers of principal components for 1D-only (diagonal) versus 1D+2D (diagonal + off-diagonal) signals (Table~\ref{tab:heritability}).

\begin{table}[H]
\centering
\caption{Heritability estimation (adjusted $R^2$) at 50 PCA components for 1D-only versus 1D+2D signals across selected phenotypes. $\Delta h^2$ is the increase from including epistatic interactions. Statistical significance assessed by paired $t$-test across phenotypes ($n = 20$).}
\label{tab:heritability}
\begin{tabular}{lcccc}
\toprule
\textbf{Phenotype} & \textbf{$h^2$ (1D only)}$\uparrow$ & \textbf{$h^2$ (1D+2D)}$\uparrow$ & \textbf{$\Delta h^2$} & \textbf{$p$-value} \\
\midrule
FT10 & $0.38$ & $0.54$ & $+0.16$ & $< 0.001$ \\
8W vernalization & $0.42$ & $0.56$ & $+0.14$ & $< 0.001$ \\
Phosphorus & $0.22$ & $0.34$ & $+0.12$ & $< 0.01$ \\
Calcium & $0.28$ & $0.39$ & $+0.11$ & $< 0.01$ \\
Potassium & $0.19$ & $0.30$ & $+0.11$ & $< 0.01$ \\
Iron & $0.24$ & $0.37$ & $+0.13$ & $< 0.01$ \\
Zinc & $0.21$ & $0.32$ & $+0.11$ & $< 0.01$ \\
Manganese & $0.26$ & $0.35$ & $+0.09$ & $< 0.05$ \\
Sodium & $0.18$ & $0.28$ & $+0.10$ & $< 0.01$ \\
Molybdenum & $0.31$ & $0.42$ & $+0.11$ & $< 0.01$ \\
\midrule
\textbf{Mean $\pm$ SD} & $\mathbf{0.27 \pm 0.08}$ & $\mathbf{0.39 \pm 0.09}$ & $\mathbf{+0.12 \pm 0.02}$ & $< 0.001$ \\
\bottomrule
\end{tabular}
\end{table}

\subsection{Machine Learning Phenotype Prediction}

Six ML models were trained on varying numbers of SNP features (50, 100, 500, 1000) derived from 1D-only or 1D+2D signals, with phenotypes discretized into 3 or 5 classes. Across all 18 ionomic phenotypes, including epistatic markers consistently improved prediction accuracy (Table~\ref{tab:ml_results}).

\begin{table}[H]
\centering
\caption{Machine learning prediction accuracy (mean $\pm$ SD across 18 phenotypes) for 1D-only versus 1D+2D SNP inputs. Accuracy measured by 5-fold cross-validation. Results shown for 500 SNP features and 3-class discretization. Statistical comparison by Wilcoxon signed-rank test across all model--phenotype combinations ($n = 108$).}
\label{tab:ml_results}
\begin{tabular}{lccccc}
\toprule
\textbf{ML Model} & \textbf{1D Only Acc.}$\uparrow$ & \textbf{1D+2D Acc.}$\uparrow$ & \textbf{$\Delta$ Acc.} & \textbf{$W$-stat} & \textbf{$p$-value} \\
\midrule
SVM \citep{cortes1995support} & $0.56 \pm 0.08$ & $0.66 \pm 0.07$ & $+0.10$ & $18$ & $< 0.001$ \\
Random Forest \citep{breiman2001random} & $0.61 \pm 0.07$ & $0.70 \pm 0.06$ & $+0.09$ & $22$ & $< 0.001$ \\
DNN \citep{lecun2015deep} & $0.58 \pm 0.09$ & $0.68 \pm 0.07$ & $+0.10$ & $16$ & $< 0.001$ \\
Gaussian Process \citep{rasmussen2006gaussian} & $0.54 \pm 0.08$ & $0.63 \pm 0.07$ & $+0.09$ & $21$ & $< 0.001$ \\
LASSO \citep{tibshirani1996regression} & $0.59 \pm 0.07$ & $0.67 \pm 0.06$ & $+0.08$ & $24$ & $< 0.001$ \\
Elastic-Net \citep{zou2005regularization} & $0.60 \pm 0.07$ & $0.68 \pm 0.06$ & $+0.08$ & $25$ & $< 0.001$ \\
\midrule
\textbf{Overall} & $\mathbf{0.58 \pm 0.07}$ & $\mathbf{0.67 \pm 0.06}$ & $\mathbf{+0.09}$ & --- & $< 10^{-8}$ \\
\bottomrule
\end{tabular}
\end{table}

The improvement from epistatic markers was robust across different numbers of input SNPs and class counts (Table~\ref{tab:ml_config}).

\begin{table}[H]
\centering
\caption{Mean prediction accuracy across all 6 ML models and 18 phenotypes for each input configuration. 1D+2D accuracy consistently exceeds 1D-only. $\Delta$ is the average improvement. Paired $t$-test across the 108 model--phenotype combinations within each configuration.}
\label{tab:ml_config}
\begin{tabular}{cccccc}
\toprule
\textbf{Classes} & \textbf{SNPs} & \textbf{1D Only}$\uparrow$ & \textbf{1D+2D}$\uparrow$ & \textbf{$\Delta$} & \textbf{$p$-value} \\
\midrule
3 & 50 & $0.51 \pm 0.08$ & $0.58 \pm 0.07$ & $+0.07$ & $< 0.001$ \\
3 & 100 & $0.54 \pm 0.08$ & $0.62 \pm 0.07$ & $+0.08$ & $< 0.001$ \\
3 & 500 & $0.58 \pm 0.07$ & $0.67 \pm 0.06$ & $+0.09$ & $< 0.001$ \\
3 & 1000 & $0.60 \pm 0.07$ & $0.69 \pm 0.06$ & $+0.09$ & $< 0.001$ \\
5 & 50 & $0.38 \pm 0.07$ & $0.44 \pm 0.06$ & $+0.06$ & $< 0.001$ \\
5 & 100 & $0.41 \pm 0.07$ & $0.48 \pm 0.06$ & $+0.07$ & $< 0.001$ \\
5 & 500 & $0.44 \pm 0.07$ & $0.52 \pm 0.06$ & $+0.08$ & $< 0.001$ \\
5 & 1000 & $0.46 \pm 0.06$ & $0.54 \pm 0.06$ & $+0.08$ & $< 0.001$ \\
\bottomrule
\end{tabular}
\end{table}

\section{Discussion}

We have demonstrated that full pairwise epistatic interaction mapping is computationally tractable using GPU acceleration and yields biologically meaningful signals that improve both heritability estimation and phenotype prediction. Three key findings emerge from this work.

First, the NGG framework reduces computation time from years to hours, making exhaustive epistatic analysis practical for the first time \citep{wei2014detecting, marchini2005genomewide}. The concordance between NGG and EMMA for 1D signals, including recovery of the well-characterized FLC flowering time locus \citep{michaels2008flowering}, validates the statistical approach. The ability to generate complete 2D epistatic maps at gene resolution enables systematic exploration of interaction architecture.

Second, epistatic interactions recover a substantial fraction of missing heritability. Across 20 phenotypes, including epistatic signals increased adjusted $R^2$ by an average of 0.12, representing approximately one-third of the gap between 1D-estimated and family-based heritability \citep{manolio2009finding, zuk2012mystery}. This provides direct empirical support for the long-standing hypothesis that epistasis contributes meaningfully to missing heritability.

Third, the improvement in ML prediction accuracy with epistatic markers is consistent across six different model architectures and 18 phenotypes. This generality suggests that the epistatic signals captured by NGG reflect genuine biological interactions rather than overfitting or model-specific artifacts. The improvement is observed even at small feature counts (50 SNPs), indicating that a modest number of high-confidence epistatic markers can substantially boost prediction.

Several limitations should be acknowledged. The MAF $> 0.3$ threshold, chosen to ensure statistical power, excludes rare variants that may also participate in important interactions. The current implementation evaluates only pairwise interactions; higher-order epistasis remains unexplored. Additionally, while we demonstrate that epistatic signals improve prediction, the biological interpretation of individual interactions requires further functional validation.

The scalability of NGG to other organisms depends on genome size and population structure. For organisms with larger genomes (e.g., human), the number of pairwise tests grows quadratically, but the GPU parallelization strategy scales favorably with additional hardware. We anticipate that the framework can be extended to human-scale datasets with modest increases in computational resources.

\section{Methods}

\subsection{Genotype and Phenotype Data}

Genotype data from the 1001 Genomes \textit{Arabidopsis thaliana} project \citep{genomes20111001} comprising 9,124,892 SNPs across 1,135 ecotypes were used. Flowering time phenotypes were from \citet{atwell2010genome}; ionomic phenotypes (18 elemental concentrations measured by ICP-MS) were from \citet{campos2021ionomics}. After MAF filtering ($> 0.3$), 341,067--371,956 SNPs were retained per phenotype analysis.

\subsection{NGG Algorithm}

For each SNP pair $(i, j)$, NGG estimates the pairwise effect $\theta_{ij}$ on the phenotype using a regression approach with the Col-0 reference genome as baseline. The diagonal elements ($\theta_{ii}$) correspond to single-SNP effects (1D-GWAS). The off-diagonal elements represent epistatic interactions. A bootstrap procedure (1000 resamples) distinguishes real signals from noise. The full computation was performed on 4 NVIDIA Quadro RTX 6000 GPUs using custom CUDA code.

\subsection{Heritability Estimation}

PCA was applied to the NGG signal matrices. Adjusted $R^2$ was computed by regressing phenotype values on increasing numbers of PCs, comparing 1D-only (diagonal elements) against 1D+2D (diagonal + upper triangle). The difference in $R^2$ between conditions was tested by paired $t$-test across phenotypes.

\subsection{Machine Learning Prediction}

Six ML models---SVM, Random Forest, DNN (2 hidden layers), Gaussian Process, LASSO, and Elastic-Net---were trained to predict discretized phenotype classes (3 or 5 classes) from SNP marker vectors (50, 100, 500, or 1000 features). Features were selected by ranking SNPs by NGG signal strength from either 1D-only or 1D+2D analyses. Accuracy was assessed by 5-fold cross-validation, stratified by phenotype class.

\subsection{Statistical Analysis}

Pearson correlation assessed concordance between NGG and EMMA signals. Paired $t$-tests and Wilcoxon signed-rank tests compared 1D-only versus 1D+2D conditions. All significance tests were two-sided with $\alpha = 0.05$. Bootstrap confidence intervals (95\%, 1000 resamples) were computed for effect size estimates.

\section{Conclusion}

Full pairwise epistatic interaction mapping is computationally tractable with GPU acceleration and retrieves biologically meaningful signals from the \textit{Arabidopsis} genome. Epistatic interactions account for a substantial fraction of missing heritability and consistently improve machine learning phenotype prediction across multiple models and phenotypes. These results establish NGG as a practical framework for exhaustive epistatic analysis and support the hypothesis that gene--gene interactions are an important component of the genetic architecture of complex traits.

\bibliographystyle{plainnat}
\bibliography{references}

\end{document}
